% Lösungen Zone – Nr. 1–12
\einheitenkopf[lzone][Kennst du schon]{Das kennst du schon}
\zweigzeile{Zuordnung: 1 → alle Einheiten · 2–4 → Einheit 1, 4 und 5 · 5 → Einheit 2 und 3 · 6 → Einheit 1 · 7 → Einheit 2 · 8 → Einheit 2 und 3 · 9, 11, 12 → Einheit 2 und 3 · 10 → Einheit 3, 4 und 5}

\erg{1}{a) 4,8 \quad b) 27 \quad c) 260 \quad d) 0,64 \quad e) 8,4 \quad f) $\approx 59{,}5$ (erst am Ende runden: 59,5125) \quad g) $\approx 4{,}2$ (4,23)}
\erg{2}{A, B, C eingetragen; C liegt mitten zwischen den Gitterlinien 2 und 3 bzw. 3 und 4}
\erg{3}{D auf der Linie 100; E zwischen 200 und 300 auf der Linie 250; F in der Mitte zwischen 350 und 400}
\erg{4}{a) 5 \quad b) 100 \quad c) 200 \quad d) 0,5}
\erg{5}{a) 330\,€ \quad b) 40\,kg \quad c) 67,20\,€ \quad d) 175,50\,€ \quad e) 36\,€ \quad f) 242\,€ (nicht 240\,€: $200 \cdot 1{,}1 \cdot 1{,}1$)}
\erg{6}{a) 14, 17 \quad b) 25, 20 \quad c) 2,1; 2,4 \quad d) 4\,€, 6\,€, 14\,€ \quad e) bei 5, im Punkt $(0\,|\,5)$}
\erg{7}{a) 3 \quad b) 2 \quad c) 1,5 \quad d) 0,9 \quad e) 1,25 \quad f) 0,8}
\erg{8}{a) 7\,€ \quad b) 20\,kg \quad c) 13,50\,€ \quad d) 15\,\% \quad e) 12\,\% (Grundwert ist der alte Preis 50\,€: $6 : 50 = 0{,}12$)}
\erg{9}{a) 1020\,€ \quad b) 1040,40\,€ \quad c) 562,43\,€ ($500 \cdot 1{,}04^3$)}
\erg{10}{a) 32 \quad b) 1000 \quad c) 3,375 \quad d) $-8$ \quad e) $\approx 0{,}66$ \quad f) $\approx 6{,}22$ (erst die Potenz, dann mal 3)}
\erg{11}{Fehler: Tom addiert die Prozentsätze; die Zinsen des Vorjahres werden aber mitverzinst. Richtig: $800\,€ \cdot 1{,}05^3 = 926{,}10\,€$}
\erg{12}{a) $1500\,€ \cdot 1{,}04^2 = 1622{,}40\,€$}
